%% ch00_orientation.tex -- Intelligence Engineering, chapter 00 "Orientation".
%%
%% The orientation lecture. It has no book chapter of its own: segments 1 to 3
%% are Huyen's preface and segment 4 is the opening of chapter 1, everything
%% before the heading "When to Use Machine Learning". ch01 opens on exactly that
%% heading, so the two files meet without overlap and without a gap.
%%
%% STATE: titles, subtitles and TEMPLATE only. No body text. Every frame carries
%% the layout it is meant to be filled into, and a % TEMPLATE comment saying why
%% that layout. The subtitles are the book's own lines, not paraphrases.
%% The earlier prose draft is kept at _archive/ch00_orientation_prose.tex.bak.
%%
%% NUMERAL. Every other chapter opens on \lecturepage{I} .. \lecturepage{XI} and
%% the numeral equals the book chapter. This lecture has no book chapter, so its
%% numeral is a blank: LECTURE _.
%%
%% IMAGES. \phimg draws a labelled placeholder box; the label is the filename to
%% drop into images/. Swapping in the asset is
%% \includegraphics[width=\linewidth]
%% in place of the \phimg call, nothing else moves.
%%
%% Fragment of intelligence_engineering.tex. One chapter is one lecture and gets
%% exactly ONE \lecturepage, at the top. Inside it every \subsection opens on
%% the subchapter divider, emitted by the master via \AtBeginSubsection.

\begin{refsection}

\lecturepage{\_}{Orientation}{what this course is, what it assumes,
  what yor are going to know after.}
\section[Orientation]{Orientation}

% ----------------------------------------------------------------------



\begin{frame}{}
  \begin{statement}
    \substatement{You have been given a lot of raw \alert{data}, }  
    and you want to engineer the shit out of it.
  \end{statement}
\end{frame}

\begin{frame}{}
  \begin{statement}
    \substatement{You have been given a business problem, }  
    and you want to choose the right \alert{metrics and objectives} to solve it.
  \end{statement}
\end{frame}

\begin{frame}{}
  \begin{statement}
    \substatement{Your initial models perform well in offline experiments, }  
    and you
    want to \alert{deploy} them at scale.
    \end{statement}
\end{frame}

\begin{frame}{}
  \begin{statement}
    \substatement{You want to harden your deployments  }  
    to quickly detect, debug, 
and \alert{address issues} that arise in production.
  \end{statement}
\end{frame}

\begin{frame}{}
  \begin{statement}
    \substatement{There is a lot of manual work and handling involved so far,  }  
    and you want to \alert{operation vacation} your machine learning product or service.
  \end{statement}
\end{frame}

\begin{frame}{}
  \begin{statement}
    \substatement{You want to learn about a framework and foundation for building machine learning systems}  
    that can be
    \alert{shared and reused} across machine learning use cases and data products.  
  \end{statement}
\end{frame}




\subchaptergloss{Machine Learning and Data Engineering}
\subsection{Overview}

% TEMPLATE: overview table. One row per lecture, the numeral equals the
% \lecturepage numeral and the title is the book's own chapter title. The
% orientation row sits above the midrule because it has no book chapter.
\begin{frame}{The lecture in one slide}
  \sideimagecover[0.33]{placeholder_gray.png}
  \sidecaption{Machine Learning Systems Overview}
  \begin{sidebody}
    \footnotesize
    \renewcommand{\arraystretch}{1.1}%
    \begin{tabular}{@{}ll@{}}
        \toprule
        Lecture & \alert{Contents}\\
        \midrule
        \alertbf{I} & Overview of Machine Learning Systems\\
        \alertbf{II} & Introduction to Machine Learning Systems Design\\
        \alertbf{III} & Data Engineering Fundamentals\\
        \alertbf{IV} & Training Data\\
        \alertbf{V} & Feature Engineering\\
        \alertbf{VI} & Model Development and Offline Evaluation\\
        \alertbf{VII} & Model Deployment and Prediction Service\\
        \alertbf{VIII} & Data Distribution Shifts and Monitoring\\
        \alertbf{IX} & Continual Learning and Test in Production\\
        \alertbf{X} & Infrastructure and Tooling for MLOps\\
        \alertbf{XI} & The Human Side of Machine Learning\\
        \bottomrule
    \end{tabular}
  \end{sidebody}
  \slidesources{Huyen2022}
\end{frame}


% ----------------------------------------------------------------------
% TEMPLATE: the one slide, unfolded. Three lectures per frame, one column
% per lecture: the chapter header on top, its subchapters below, titles
% only. The frame title stays the overview frame's. The size command sits
% inside each column on purpose: a bare group wrapping the columns block
% silently drops it.
\begin{frame}{The lecture in one slide}
  \begin{columns}[t]
    \begin{column}{0.31\textwidth}
      \footnotesize
      \renewcommand{\arraystretch}{1.1}%
      \begin{tabular}[t]{@{}>{\raggedright\arraybackslash}p{\linewidth}@{}}
        \toprule
        \parbox[t][2\baselineskip][t]{\linewidth}{\raggedright\alertbf{I}\quad \textbf{Overview of Machine Learning Systems}}\\
        \midrule
        When to Use Machine Learning\\[0.2em]
        Understanding Machine Learning Systems\\[0.2em]
        \bottomrule
      \end{tabular}
    \end{column}
    \begin{column}{0.31\textwidth}
      \footnotesize
      \renewcommand{\arraystretch}{1.1}%
      \begin{tabular}[t]{@{}>{\raggedright\arraybackslash}p{\linewidth}@{}}
        \toprule
        \parbox[t][2\baselineskip][t]{\linewidth}{\raggedright\alertbf{II}\quad \textbf{Introduction to Machine Learning Systems Design}}\\
        \midrule
        Business and ML Objectives\\[0.2em]
        Requirements for ML Systems\\[0.2em]
        Iterative Process\\[0.2em]
        Framing ML Problems\\[0.2em]
        Mind Versus Data\\[0.2em]
        \bottomrule
      \end{tabular}
    \end{column}
    \begin{column}{0.31\textwidth}
      \footnotesize
      \renewcommand{\arraystretch}{1.1}%
      \begin{tabular}[t]{@{}>{\raggedright\arraybackslash}p{\linewidth}@{}}
        \toprule
        \parbox[t][2\baselineskip][t]{\linewidth}{\raggedright\alertbf{III}\quad \textbf{Data Engineering Fundamentals}}\\
        \midrule
        Data Sources\\[0.2em]
        Data Formats\\[0.2em]
        Data Models\\[0.2em]
        Data Storage Engines and Processing\\[0.2em]
        Modes of Dataflow\\[0.2em]
        Batch Processing Versus Stream Processing\\[0.2em]
        \bottomrule
      \end{tabular}
    \end{column}
  \end{columns}
  \slidesources{Huyen2022}
\end{frame}

\begin{frame}{The lecture in one slide}
  \begin{columns}[t]
    \begin{column}{0.31\textwidth}
      \footnotesize
      \renewcommand{\arraystretch}{1.1}%
      \begin{tabular}[t]{@{}>{\raggedright\arraybackslash}p{\linewidth}@{}}
        \toprule
        \parbox[t][2\baselineskip][t]{\linewidth}{\raggedright\alertbf{IV}\quad \textbf{Training Data}}\\
        \midrule
        Sampling\\[0.2em]
        Labeling\\[0.2em]
        Class Imbalance\\[0.2em]
        Data Augmentation\\[0.2em]
        \bottomrule
      \end{tabular}
    \end{column}
    \begin{column}{0.31\textwidth}
      \footnotesize
      \renewcommand{\arraystretch}{1.1}%
      \begin{tabular}[t]{@{}>{\raggedright\arraybackslash}p{\linewidth}@{}}
        \toprule
        \parbox[t][2\baselineskip][t]{\linewidth}{\raggedright\alertbf{V}\quad \textbf{Feature Engineering}}\\
        \midrule
        Learned Features Versus Engineered Features\\[0.2em]
        Common Operations\\[0.2em]
        Data Leakage\\[0.2em]
        Engineering Good Features\\[0.2em]
        \bottomrule
      \end{tabular}
    \end{column}
    \begin{column}{0.31\textwidth}
      \footnotesize
      \renewcommand{\arraystretch}{1.1}%
      \begin{tabular}[t]{@{}>{\raggedright\arraybackslash}p{\linewidth}@{}}
        \toprule
        \parbox[t][2\baselineskip][t]{\linewidth}{\raggedright\alertbf{VI}\quad \textbf{Model Development and Offline Evaluation}}\\
        \midrule
        Model Development and Training\\[0.2em]
        Model Offline Evaluation\\[0.2em]
        \bottomrule
      \end{tabular}
    \end{column}
  \end{columns}
  \slidesources{Huyen2022}
\end{frame}

\begin{frame}{The lecture in one slide}
  \begin{columns}[t]
    \begin{column}{0.31\textwidth}
      \footnotesize
      \renewcommand{\arraystretch}{1.1}%
      \begin{tabular}[t]{@{}>{\raggedright\arraybackslash}p{\linewidth}@{}}
        \toprule
        \parbox[t][2\baselineskip][t]{\linewidth}{\raggedright\alertbf{VII}\quad \textbf{Model Deployment and Prediction Service}}\\
        \midrule
        ML Deployment Myths\\[0.2em]
        Batch Prediction Versus Online Prediction\\[0.2em]
        Model Compression\\[0.2em]
        ML on the Cloud and on the Edge\\[0.2em]
        \bottomrule
      \end{tabular}
    \end{column}
    \begin{column}{0.31\textwidth}
      \footnotesize
      \renewcommand{\arraystretch}{1.1}%
      \begin{tabular}[t]{@{}>{\raggedright\arraybackslash}p{\linewidth}@{}}
        \toprule
        \parbox[t][2\baselineskip][t]{\linewidth}{\raggedright\alertbf{VIII}\quad \textbf{Data Distribution Shifts and Monitoring}}\\
        \midrule
        Causes of ML System Failures\\[0.2em]
        Data Distribution Shifts\\[0.2em]
        Monitoring and Observability\\[0.2em]
        \bottomrule
      \end{tabular}
    \end{column}
    \begin{column}{0.31\textwidth}
      \footnotesize
      \renewcommand{\arraystretch}{1.1}%
      \begin{tabular}[t]{@{}>{\raggedright\arraybackslash}p{\linewidth}@{}}
        \toprule
        \parbox[t][2\baselineskip][t]{\linewidth}{\raggedright\alertbf{IX}\quad \textbf{Continual Learning and Test in Production}}\\
        \midrule
        Continual Learning\\[0.2em]
        Test in Production\\[0.2em]
        \bottomrule
      \end{tabular}
    \end{column}
  \end{columns}
  \slidesources{Huyen2022}
\end{frame}

\begin{frame}{The lecture in one slide}
  \begin{columns}[t]
    \begin{column}{0.31\textwidth}
      \footnotesize
      \renewcommand{\arraystretch}{1.1}%
      \begin{tabular}[t]{@{}>{\raggedright\arraybackslash}p{\linewidth}@{}}
        \toprule
        \parbox[t][2\baselineskip][t]{\linewidth}{\raggedright\alertbf{X}\quad \textbf{Infrastructure and Tooling for MLOps}}\\
        \midrule
        Storage and Compute\\[0.2em]
        Development Environment\\[0.2em]
        Resource Management\\[0.2em]
        ML Platform\\[0.2em]
        Build Versus Buy\\[0.2em]
        \bottomrule
      \end{tabular}
    \end{column}
    \begin{column}{0.31\textwidth}
      \footnotesize
      \renewcommand{\arraystretch}{1.1}%
      \begin{tabular}[t]{@{}>{\raggedright\arraybackslash}p{\linewidth}@{}}
        \toprule
        \parbox[t][2\baselineskip][t]{\linewidth}{\raggedright\alertbf{XI}\quad \textbf{The Human Side of Machine Learning}}\\
        \midrule
        User Experience\\[0.2em]
        Team Structure\\[0.2em]
        Responsible AI\\[0.2em]
        \bottomrule
      \end{tabular}
    \end{column}
    \begin{column}{0.31\textwidth}
    \end{column}
  \end{columns}
  \slidesources{Huyen2022}
\end{frame}



% ----------------------------------------------------------------------
% Course work, two frames under one title like the repeated "lecture in
% one slide" frames. First the contract: format in the block title, the
% project brief in its body, logistics in grey at the bottom. The judging
% lens gets the second frame. Side strip on both: the bee hive.
\begin{frame}{Course Work}
  \sideimagecover[0.3]{bees}
  \sidecaption{The bee}
  \begin{sidebody}
    \begin{block}{Presentation (15\,min) + Q\&A (open)}
      You get recordings of a bee hive entrance, and you are to
      contribute by data engineering, an annotation policy, a prediction
      service or specialized tooling.
    \end{block}
    \vskip0.8em
    {\color{tddim}Open whiteboard, three slides or one page only.
    Be ready when the class starts.}
  \end{sidebody}
\end{frame}

% The judging lens: bold term + grey rest, one row per talk section, in
% the order a talk runs.
\begin{frame}{Course Work}
  \sideimagecover[0.3]{bee_mloverlay}
  \sidecaption{The bee}
  \begin{sidebody}
    What I am looking for ..\\
    \medskip

    \textbf{\alertbf{I}ntroduction}\enskip{\color{tddim}by Motivation and Literature}\\[0.5em]
    \textbf{\alertbf{M}ethod}\enskip{\color{tddim}by Explanation and or Demonstration}\\[0.5em]
    \textbf{\alertbf{R}esults}\enskip{\color{tddim}show what was achieved and how well}\\[0.5em]
    \textbf{\alertbf{D}iscussion}\enskip{\color{tddim}on findings, limitations and what's next}\\[0.5em]
    \textbf{Smart and Aesthetic}\enskip{\color{tddim}Thoughts, Designs, Ideas and \alert{Questions}.}
  \end{sidebody}
\end{frame}

% ----------------------------------------------------------------------
% The extra mile. The offer in the one block, the terms in grey under it,
% and the trade as the closing line carrying the slide's single red peak.
% Side strip: the large banana picture.
\begin{frame}{Extra Mile}
  \sideimagecover[0.33]{banana_1}
  \sidecaption{A large banana}
  \begin{sidebody}
    \begin{block}{Special projects}
      If you are in for an extra mile, let me know. I have special
      projects that I can assign to you.
    \end{block}
    \vskip0.8em
    {\color{tddim}Wrapped in the same presentation and Q\&A format, yes it will be about Bananas, and yes it will be
    more challenging, more time and effort.}

    \bigskip
    Classic loose loose short term, yet \alertbf{win long term}.
  \end{sidebody}
\end{frame}


\begin{frame}{}
  \slideqr{The Hive is Real. Dataset on Hugging Face.}{https://huggingface.co/datasets/mrkschtr/thehiveisreal}
  \begin{statement}
    \substatement{By the end of this lecture, }
    Get yourself a \alert{presentation slot} and brainstorm a project \alert{idea to work} on.
    \end{statement}
\end{frame}

% ----------------------------------------------------------------------
% This chapter's own references. The refsection opened above scopes the
% citation numbering to this chapter, so chapter_pdfs/<this>.pdf resolves
% its own [n] markers. Two columns and a step down in size: this lecture
% cites the whole bib and it does not fit one column at \footnotesize.
\begin{frame}{References}
  \renewcommand*{\bibfont}{\scriptsize}%
  \setlength{\bibitemsep}{0.25em}%
  \begin{multicols}{2}
    \printbibliography[heading=none]
  \end{multicols}
\end{frame}
\end{refsection}
